%************************************************************************* %
%**************************  MODELO TCC UEL  ***************************** %
%************************************************************************* %

%****************************** CLASSE *********************************** %

\documentclass[12pt, openright, twoside, a4paper, english, brazil]{abntex2}   % Padrão Abntex2
		 % openright: capítulos começam em pág ímpar
		 % twoside: para impressão em verso e anverso.

%***************************** PACOTES *********************************** %
%\usepackage[within=none]{newfloat}
\usepackage[utf8]{inputenc}		     % Codificacao do documento (conversão automática dos acentos)

%	DEEL
\usepackage{DEEL}               	 % personalização da ABNTEX2 para a Universidade Estadual de Londrina

%	Fonte
\usepackage[T1]{fontenc}		   	 % Selecao de codigos de fonte.
\usepackage[brazil]{babel}
\usepackage{lmodern}			     % Usa a fonte Latin Modern

%	Matemático e Gráfico
\usepackage{float}
\usepackage{graphics,graphicx}	     % pacotes para inserir figuras .eps ou .jpg
\usepackage{amssymb}           		 % pacote para fontes e simbolos matemáticos
\usepackage{longtable}          	 % possibilita inserir grandes tabelas
\usepackage{color,colortbl,multirow} % permite textos e tabelas com cores
\usepackage{amsmath}           		 % pacote para equações
\usepackage{epstopdf}
\usepackage{ragged2e}
\usepackage{listings}
%\lstset{numbers=left, numberstyle=\tiny, stepnumber=1, numbersep=5pt, basicstyle=\scriptsize ,frame=tbrl}
\usepackage{etoolbox}
%\usepackage{caption}
\usepackage[position=b,singlelinecheck=on]{subfig}

\floatstyle{plaintop}
\newfloat{quadro}{htbp}{loq}
\floatname{quadro}{Quadro}

\makeatletter
\patchcmd{\listof}% <cmd>
{\float@listhead}% <search>
{\@namedef{l@#1}{\l@table}\float@listhead}% <replace>
{}{}% <success><failure>
\makeatother

%\makeatletter
%\renewcommand*{\float@listhead}[1]{%
%	\@ifundefined{chapter}{%
%		\section*{#1}%
%		\addcontentsline{toc}{section}{#1}%
%	}{%
%	\chapter*{#1}% 
%	\addcontentsline{toc}{chapter}{#1}%
%}%
%\@mkboth{\MakeUppercase{#1}}{\MakeUppercase{#1}}%
%}
%\makeatother




%\DeclareFloatingEnvironment{quadro}[Quadro][Lista de quadros]
%\DeclareCaptionType{quadro}[Quadros][Lista de quadros]


% % % % % % % % % % % % % % % % % % % %

\definecolor{codegreen}{rgb}{0,0.6,0}
\definecolor{codegray}{rgb}{0.5,0.5,0.5}
\definecolor{codepurple}{rgb}{0.58,0,0.82}
\definecolor{backcolour}{rgb}{0.95,0.95,0.92}

\definecolor{verde}{rgb}{0,.5725,0.247}

\lstdefinestyle{mystyle}{
	backgroundcolor=\color{backcolour},   
	commentstyle=\color{codegreen},
	keywordstyle=\color{magenta},
	numberstyle=\tiny\color{codegray},
	stringstyle=\color{codepurple},
	basicstyle=\footnotesize,
	breakatwhitespace=false,         
	breaklines=true,                 
	captionpos=b,                    
	keepspaces=true,                 
	numbers=left,                    
	numbersep=5pt,                  
	showspaces=false,                
	showstringspaces=false,
	showtabs=false,                  
	tabsize=2
}

\lstset{style=mystyle}

%	Estrutural
\usepackage{url}
\usepackage{threeparttable}     	 % permite a inserção de notas de rodapé nas tabelas
\usepackage{lscape}             	 % orientação de página LANDSCAPE
\usepackage{pifont}             	 % acrescenta símbolos diferentes.
%\usepackage{cmap}					 % Mapear caracteres especiais no PDF	
\usepackage{lastpage}				 % Usado pela Ficha catalográfica
\usepackage{indentfirst}			 % Indenta o primeiro parágrafo de cada sessão.



% Pacotes de citações
\usepackage[brazilian,hyperpageref]{backref}	 % Paginas com as citações na bibl
\usepackage[alf]{abntex2cite}	% Citações padrão ABNT

\data{2016}        % Ano de realização

\begin{document}

%---------------------------------------------------------------
%--------- DADOS BÁSICOS DO TRABALHO (Preencher todos) ---------
%---------------------------------------------------------------


\titulo{Não vai ter golpe}   % Título do Trabalho em Português

\tituloingles{There won't be coup}   % Título do Trabalho em Inglês

\palavraschave{1. Tecnologia Assistiva. 2. Tetraplegia. 3. Transdutores. 4. Baixo Custo.}  % 5 palavras-chave
\keywords{1. Assistive Technology. 2. Tetraplegia. 3. Transducers. 4. Low Cost.}   % 5 key-words

\DEELautor{Dilma}{Vana Rousseff}                  % Aluno autor do trabalho

\curso{Engenharia Elétrica}												% Ex: Engenharia Civil

\orientador{Prof. Dr. Luiz Inácio Lula da Silva}              % Professor Orientador membro da banca 

\membrob{Prof. Dr. Maria do Rosário}   % Professor membro da banca 
\membroc{Prof. Dr. Renan Calheiros}   % Professor membro da banca

\local{Londrina}   % Cidade da instituição

%--------------------------------------------------------------------------------------------


\DEELdedicatoria{Dedico este trabalho a todos aqueles que, de alguma forma,\\ auxiliaram para a concretização desta etapa. \\ }  % Texto da Dedicatória (OBS: dedicatória é opcional)
                                    % Caso não queira inserir, deixe em branco ( \dedicatoria{} )

\DEELagradecimentos{Agradeço primeiramente a CUT e ao MST, por fornecer-me as condições necessárias para conclusão desta importante etapa da minha vida.} 
                                   % Texto dos Agradecimentos (OBS: Agradecimentos é opcional)
                                   % Caso não queira inserir, deixe em branco ( \agradecimentos{} )
                                  

\DEELepigrafe{\textit{"O dia da criança é dia da mãe,\\ do pai e das professoras,\\ mas também é o dia dos animais.\\ Sempre que você olha uma criança, há sempre\\ uma figura oculta, que é um cachorro atrás,\\ o que é algo muito importante."} \\ (Dilma Vana Rousseff)} 
                                   % Texto da Epígrafe (OBS: Epígrafe é opcional)
                                   % Caso não queira inserir, deixe em branco ( \epigrafe{} )


\DEELresumo{Blablablabla yada yada blabla bla} 
                                   % Texto do Resumo (OBS: Resumo é Obrigatório)

\DEELabstract{Blablablabla yada yada blabla bla} 
                                   % Texto do Abstract (OBS: Abstract é Obrigatório)



%----- Elementos Pré-Textuais--------------------------
\ElementosPreTextuais
% Insere todos os elementos Pré-Textuais listados (capa, folha de rosto, ...)
												% capa                   % Obrigatório
												% folhaderosto           % Obrigatório
												% fichacatalografica     % Obrigatório
												% folhaaprovacao         % Obrigatório
												% dedicatoria            % Opcional
												% agradecimentos         % Opcional
												% epigrafe               % Opcional
												% resumo                 % Obrigatório
												% abstract               % Obrigatório
\
%--------------------------------------------------------------------------------------------
%\newpage
%---- Outros Elementos Pré-Textuais ----
%--------Lista de figuras--------------------------------------------------------------------
\listoffigures*
\cleardoublepage
%--------Lista de tabelas--------------------------------------------------------------------
\listoftables*
\cleardoublepage

\renewcommand{\cfttablename}{Quadro\space} 
\listof{quadro}{Lista de quadros}
\cleardoublepage

%--------Lista de Siglas e Abreviaturas------------------------------------------------------
\include{list_siglas}   % (OBS: Elemento Condicionado à Necessidade) 
%                       % Veja modelo no arquivo list_siglas.tex
%--------Lista de Simbolos e Notação --------------------------------------------------------
%\include{list_simbolos} % (OBS: Elemento Condicionado à Necessidade)
%                       % Veja modelo no arquivo list_simbolos.tex
%--------Sumário-----------------------------------------------------------------------------
\cleardoublepage
\tableofcontents*
\thispagestyle{empty}
%--------------------------------------------------------------------------------------------


%----- Elementos Textuais--------------------------
\textual 
\include{introducao}				  % Capítulo 1 - Texto do cap1.tex
%%
\include{fundamentacao}				  % Capítulo 2 - Texto do cap2.tex
%%
\include{desenvolvimento}			  % Capítulo 3 - Texto do cap3.tex
%%
\include{resultados}		          % Capítulo 4 - Texto do cap4.tex
%%
\include{conclusao}				      % Capítulo 5 - Texto do cap5.tex
%%
%% Insira ou retire quantos capítulos desejar
%%--------------------------------------------------------------------------------------------


%----- Elementos Pós-Textuais--------------------------
%
%--------Referências-------------------------------------------------------------------------
\postextual
\bibliography{bibliografia_teste} % arquivo bibtex com as referências (modelo.bib)

%%--------Apendice-------------------------------------------------
%\apendice                  % Comando que inclui os apendices (sua autoria)
%%
%\include{apendice_1}				  % Apendice 1 - arquivo apendice_1.tex
%%
%\include{apendice_2}         % Apendice 2 - arquivo apendice_2.tex
%%
%% Descomente caso exista a necessidade da inclusão de apendices
%%---------------------------------------------------------------

%%--------Anexos-------------------------------------------------
%\anexo                     % Comando que inclui os anexos (outros autores)
%%
%\include{anexo_1}				  % Anexo 1 - arquivo anexo_1.tex
%%
%\include{anexo_2}          % Anexo 2 - arquivo anexo_2.tex
%%
%% Descomente caso exista a necessidade da inclusão de anexos
%%---------------------------------------------------------------

\end{document}
